% =============================================================================
% EBF WORKING PAPER WP-001
% =============================================================================
%
% Format:     Working Paper
% Seiten:     ~30
% Qualität:   Q2
% Ziel:       Mehrdimensionales Framework für Torwarttechnik-Bewertung
% Leser:      Sportwissenschaft, Torwarttrainer-Community, Coaching-Forschung
%
% Nummerierung: WP-001
% Versionierung: v3
%
% =============================================================================

\documentclass[12pt,a4paper]{article}
\usepackage[margin=2.5cm]{geometry}
\usepackage[utf8]{inputenc}
\usepackage[T1]{fontenc}
\usepackage[ngerman]{babel}
\usepackage{amsmath,amssymb}
\usepackage{booktabs}
\usepackage{tcolorbox}
\usepackage{hyperref}
\usepackage{natbib}
\usepackage{array}
\usepackage{tabularx}
\usepackage{parskip}
\usepackage{graphicx}
\usepackage{xcolor}
\usepackage{float}
\usepackage{enumitem}

% =============================================================================
% FEHRADVICE FARBEN
% =============================================================================
\definecolor{fadarkblue}{RGB}{2,64,121}
\definecolor{falightblue}{RGB}{84,158,222}
\definecolor{fadarkgray}{RGB}{37,33,42}
\definecolor{falightgray}{RGB}{243,245,247}

% =============================================================================
% HYPERLINK SETUP
% =============================================================================
\hypersetup{
  colorlinks=true,
  linkcolor=fadarkblue,
  citecolor=fadarkblue,
  urlcolor=fadarkblue
}

% =============================================================================
% METADATEN
% =============================================================================
\newcommand{\WPNumber}{WP-001}
\newcommand{\WPVersion}{v3}
\newcommand{\WPDate}{Februar 2026}
\newcommand{\WPTitle}{Torwart-Leistung im Leistungsverbund: Ein mehrdimensionales Framework zur Bewertung von Torwarttechniken und Trainingsallokation -- mit Visibility Framework, Counterfactual Performance, Action Bias, Verletzungsrisiko und Verf\"ugbarkeit}
\newcommand{\WPAuthors}{Gerhard Bruno Fehr}
\newcommand{\WPStatus}{Draft}

% =============================================================================
\begin{document}
% =============================================================================

% -----------------------------------------------------------------------------
% TITELSEITE
% -----------------------------------------------------------------------------
\begin{titlepage}
\begin{center}

\vspace*{1cm}

{\small\textsc{Working Paper}}\\[1.5cm]

\rule{\linewidth}{0.5pt}\\[0.4cm]
{\LARGE\bfseries \WPTitle}\\[0.4cm]
\rule{\linewidth}{0.5pt}\\[2cm]

{\Large \WPAuthors}\\[0.3cm]
{\large FehrAdvice \& Partners AG, Z\"urich}\\[2.5cm]

\begin{tcolorbox}[colback=fadarkblue!5!white, colframe=fadarkblue, width=8cm]
\centering
\textbf{Working Paper \WPNumber}\\
Version \WPVersion\\
\WPDate\\[0.5em]
Status: \textbf{\WPStatus}
\end{tcolorbox}

\vfill

{\small
\textbf{Abstract}\\[0.5em]
\begin{minipage}{0.85\textwidth}
Die aktuelle Torwarttrainer-Praxis steht unter dem Einfluss eines Social-Media-getriebenen Trends zur \"Uberbetonung von Blocktechniken (Spread, K-Block, Langblock). Dieses Working Paper entwickelt ein mehrdimensionales Framework zur systematischen Bewertung von Torwarttechniken und deren Allokation in der Trainingsplanung. Das Modell unterscheidet drei Leistungsebenen -- strategisch, taktisch, operativ -- und verortet jede Ebene im Leistungsverbund mit der Mannschaft. Techniken werden anhand von \textbf{sechs} Kriterien bewertet: Spielrelevanz, Risikoreduktion, Lernaufwand, Trainerkompetenz, strategisches Leistungspotential und \textbf{Verletzungsrisiko} (differenziert nach Verletzungen ohne und mit Fremdeinwirkung). Die Analyse zeigt, dass Fangen in f\"unf von sechs Dimensionen \"uberlegen ist -- Blocken verliert auf \emph{allen sechs} Dimensionen und birgt zus\"atzlich das h\"ochste Verletzungsrisiko aller Techniken. Die Validierung durch das Red-Bull-Leipzig/Salzburg-Scouting-Framework (Wagner, 2023) best\"atigt die Praxisrelevanz. Das Paper schl\"agt ein datenbasiertes Trainingsprotokoll vor und f\"uhrt die Verf\"ugbarkeitsformel $V^{*}_{\text{eff}} = V^{*} \times A \times D$ ein: Selbst die beste Technik ist wertlos, wenn der Torh\"uter verletzt ist.
\end{minipage}
}

\vspace{0.8cm}

{\footnotesize
Keywords: Torwarttraining, Blocktechnik, Fangen, Trainingsallokation, Leistungsverbund,\\
Komplementarit\"at, Entscheidungsqualit\"at, Verletzungsrisiko, Verf\"ugbarkeit, Red Bull Scouting
}

\end{center}
\end{titlepage}

% -----------------------------------------------------------------------------
% VERSION HISTORY
% -----------------------------------------------------------------------------
\begin{tcolorbox}[colback=gray!5!white, colframe=gray!50!black,
    title=\textbf{Version History}]
\begin{tabular}{@{}llp{8cm}@{}}
\toprule
\textbf{Version} & \textbf{Datum} & \textbf{\"Anderungen} \\
\midrule
v1 & Februar 2026 & Initial draft -- Diskussionsentwurf \\
v2 & Februar 2026 & Literatur-Integration: Visibility Framework (Holmstr\"om-Milgrom), Counterfactual Performance (Rose), Action Bias (Bar-Eli), Positions-Generalisierung \\
v3 & Februar 2026 & 6.\ Dimension Verletzungsrisiko (ohne/mit Fremdeinwirkung), Red-Bull-Scouting-Validierung (Wagner 2023), Verf\"ugbarkeitsformel $V^{*}_{\text{eff}} = V^{*} \times A \times D$, Verletzungs-Substitution als 4.\ Substitutionskostenkategorie \\
\bottomrule
\end{tabular}
\end{tcolorbox}

\vspace{0.5cm}
{\small\textit{Dieses Working Paper ist ein Diskussionsentwurf und wird laufend weiterentwickelt.}}

\tableofcontents
\newpage

% =============================================================================
\section{Einleitung: Die Block-Debatte}
\label{sec:einleitung}
% =============================================================================

In den vergangenen Jahren hat sich in der Torwarttraining-Community ein deutlicher Trend zur Betonung von Blocktechniken entwickelt. Spread-Saves, K-Blocks und Langblocks werden auf Social-Media-Plattformen als zentrale Elemente modernen Torwartspiels gefeiert. Virale Clips spektakul\"arer Block-Saves erzeugen den Eindruck, diese Technik sei die L\"osung schlechthin f\"ur das moderne Torwartspiel.

Gleichzeitig fehlt es an einer systematischen Grundlage, um den tats\"achlichen Wert einzelner Techniken im Verh\"altnis zur gesamten Torwart-Leistung zu bewerten. Die klassischen Metriken -- Save Percentage, Expected Goals (xG), Expected Goals Against (xGA) -- erfassen die operative Ergebnisqualit\"at, ignorieren jedoch die taktischen und strategischen Entscheidungen, die der Aktion vorausgehen und ihr folgen. In der klassischen Bewertung gehen gerade die taktischen und strategischen Dimensionen der Torwartleistung komplett verloren, obwohl sie den Spielverlauf massiv beeinflussen k\"onnen \citep{tigges2026}.

Dieses Paper stellt die zentrale Frage: \textbf{Nach welchen Kriterien sollte Trainingszeit auf verschiedene Torwarttechniken allokiert werden? Und welches Leistungsmodell bildet die Grundlage f\"ur diese Allokation?}

Die Kernthese lautet: Die Funktion von Trainingszeit l\"asst sich auf drei Prinzipien reduzieren -- die H\"aufigkeitsverteilung des Spiels abbilden, die Risikoreduktion maximieren und das strategische Leistungspotential des Torh\"uters im Leistungsverbund mit der Mannschaft entwickeln.


% =============================================================================
\section{Empirische Ausgangslage}
\label{sec:empirie}
% =============================================================================

\subsection{Verf\"ugbare Daten}
\label{sec:daten}

Die Datenlage zur spezifischen Aufschl\"usselung von Save-Typen ist \"uberraschend d\"unn. Opta erfasst Save-Typen granular (Caught, Collected, Parried Safe, Parried Danger, Fumble, Fingertip), kombiniert mit K\"orperteil (Hands, Feet, Body, Fist) und Bewegungsmuster (Diving, Standing, Sliding, Reaching, Stooping). Diese Rohdaten werden jedoch nicht aggregiert \"offentlich zug\"anglich gemacht \citep{opta2024}.

Verf\"ugbare Referenzwerte: Torh\"uter in den Top-5-Ligen parieren im Schnitt 69\% aller Sch\"usse aufs Tor. Pro Spiel stehen durchschnittlich 4--5 Sch\"usse aufs Tor pro Torh\"uter an \citep{soccerment2024}. Die K\"orperfl\"ache eines durchschnittlichen Top-Liga-Torh\"uters (190\,cm) deckt stehend bereits 24\% der Torfl\"ache ab; die verbleibenden 45 Prozentpunkte der Save-Quote resultieren aus Positionierung, Agilit\"at und Reflexen.


\subsection{Die Stichprobenproblematik: Noise vs.\ Signal}
\label{sec:noise}

Bei 4--5 Sch\"ussen aufs Tor pro Spiel ist die Stichprobe f\"ur jede einzelne Save-Technik extrem klein. Ein einzelner spektakul\"arer Block-Save kann reines statistisches Rauschen sein. Wird dieses Ereignis auf Social Media verbreitet, entsteht ein klassischer Verf\"ugbarkeitsbias: Noise wird als Signal interpretiert, und eine Einzelsituation wird zur Trainingsphilosophie erhoben.

Die Studie von \citet{gramage2025} befragte \"uber 400 Torwarttrainer international und zeigt: 99\% der Trainer erkennen die Spread-Technik, 78\% setzen sie aktiv im Training ein, mit einer empfohlenen Einf\"uhrung ab ca.\ 11 Jahren. Dieser hohe Trainingsanteil steht in einem bemerkenswerten Missverh\"altnis zur tats\"achlichen H\"aufigkeit reiner 1v1-Blocksituationen im Spiel.


\subsection{Was xG und xGA nicht messen}
\label{sec:xg-limitationen}

Die g\"angigen Torwart-Metriken (xG, xGA, Post-Shot xG, GSAA) messen die Qualit\"at der Schuss-Situation und die Save-Wahrscheinlichkeit. Was sie \emph{nicht} erfassen:

\begin{itemize}[nosep]
  \item Hat der Torh\"uter durch \textbf{Positionierung} die Schussqualit\"at bereits reduziert, \emph{bevor} der Schuss fiel?
  \item Hat er durch \textbf{Coaching der Abwehr} die Situation entsch\"arft?
  \item Was passiert \textbf{nach dem Save} -- Ballbesitz gesichert und Tempogegenstoss eingeleitet, oder Ball im Spiel und R\"uckkehr in die Defensive?
\end{itemize}


% =============================================================================
\section{Torwart-Leistungsmodell: Drei Ebenen im Verbund}
\label{sec:leistungsmodell}
% =============================================================================

Wir schlagen ein Leistungsmodell vor, das die Torwartleistung in drei interdependente Ebenen zerlegt und jede Ebene im Leistungsverbund mit der Mannschaft verortet.


\subsection{Strategische Leistung}
\label{sec:strategisch}

Die strategische Leistung beschreibt den Beitrag des Torh\"uters zum Spielkonzept der Mannschaft: Spieler\"offnung, Tempokontrolle, Raumverkn\"appung, \"Uberzahl im Aufbauspiel. Die zentrale Frage: \emph{Macht der Torh\"uter das Team als Ganzes besser?}

\textbf{Strategischer Leistungsverbund:} Das Spielkonzept des Trainers definiert, was vom Torh\"uter strategisch gefordert wird. Hohes Pressing erfordert Absicherung und Spieler\"offnung; tiefes Verteidigen ver\"andert das Anforderungsprofil grundlegend. Das ist keine Einbahnstrasse.


\subsection{Taktische Leistung}
\label{sec:taktisch}

Die taktische Leistung beschreibt die situative Entscheidungsqualit\"at in Echtzeit. Wann kommt der Torh\"uter heraus, wann bleibt er auf der Linie? Wann f\"angt er, wann pariert er, wann blockt er?

\textbf{Taktischer Leistungsverbund:} Die Entscheidungen des Torh\"uters h\"angen direkt von der Abwehrkette ab und umgekehrt. Kommunikation, Timing und gegenseitiges Lesen sind taktische Ko-Produktion. Ein blockender Torh\"uter zwingt seine Mannschaft zur\"uck in die Defensivarbeit; ein fangender Torh\"uter entlastet das System.


\subsection{Operative Leistung}
\label{sec:operativ}

Die operative Leistung beschreibt die reine technische Ausf\"uhrungsqualit\"at: Fangen, Parieren, Blocken, Fausten, Fussabwehr.

\textbf{Operativer Leistungsverbund:} Die technische Qualit\"at bestimmt direkt die Anschlussaktion der Mitspieler. Ein gefangener Ball erm\"oglicht sofortiges Umschalten. Ein abgeblockter Ball erzeugt eine unklare Situation f\"ur die gesamte Mannschaft.


\subsection{Hochgradige Abh\"angigkeiten}
\label{sec:abhaengigkeiten}

Die drei Leistungsebenen stehen in einem Verh\"altnis hochgradiger Abh\"angigkeit. Die Gesamtleistung des Torh\"uters ist eine Funktion aller drei Ebenen im Mannschaftskontext:

\begin{equation}
\label{eq:leistung}
L(\text{TH}) \;=\; f\!\left(\, L_{\text{strategisch}},\; L_{\text{taktisch}},\; L_{\text{operativ}} \,\right) \;\Big|\; \text{Leistungsverbund Mannschaft}
\end{equation}


% =============================================================================
\section{Sechs-Dimensionen-Bewertungsframework}
\label{sec:framework}
% =============================================================================

Zur systematischen Bewertung jeder Torwarttechnik schlagen wir \textbf{sechs} Dimensionen vor. Die sechste Dimension -- Verletzungsrisiko -- wurde in v3.0 hinzugef\"ugt, nachdem die Analyse des Red-Bull-Scouting-Frameworks (Abschnitt~\ref{sec:redbull}) zeigte, dass Praktiker die Robustheit eines Torh\"uters als eigenst\"andige Kompetenz bewerten.

\begin{table}[H]
\centering
\caption{Sechs Bewertungsdimensionen f\"ur Torwarttechniken (v3.0)}
\label{tab:dimensionen}
\begin{tabularx}{\textwidth}{@{}l l X X@{}}
\toprule
\textbf{Dim.} & \textbf{Gew.} & \textbf{Definition} & \textbf{Messfrage} \\
\midrule
D1: Spielrelevanz & 25\% & Wie h\"aufig kommt die Situation im Spiel vor? & Anteil an Gesamtaktionen/Spiel \\[4pt]
D2: Risikoreduktion & 22\% & Senkt die Technik die Torwahrscheinlichkeit? & Ball gesichert vs.\ Ball im Spiel \\[4pt]
D3: Lernaufwand & 13\% & Trainingszeit bis zur Abrufbarkeit unter Druck? & Stunden bis Automatisierung \\[4pt]
D4: Trainerkompetenz & 8\% & Kann der Trainer situationsgerecht vermitteln? & Diagnostik-/Designf\"ahigkeit \\[4pt]
D5: Strateg.\ Pot. & 17\% & Wertsch\"opfung \"uber Gefahrenabwehr hinaus? & Transition / Spieler\"offnung \\[4pt]
\textbf{D6: Verletzungsrisiko} & \textbf{15\%} & \textbf{Wie hoch ist das Verletzungsrisiko?} & \textbf{D6a + D6b (s.u.)} \\
\bottomrule
\end{tabularx}
\end{table}

\textbf{Dimension D6} setzt sich aus zwei gleichgewichteten Subkomponenten zusammen:

\begin{itemize}[nosep]
  \item \textbf{D6a -- Verletzungsrisiko ohne Fremdeinwirkung:} Verletzungen durch die Technik selbst -- biomechanischer Stress, extreme Gelenkwinkel, repetitive Belastung. Direkt durch Technikwahl und Trainingsdesign beeinflussbar. Beispiel: Die Spread-Technik erfordert eine H\"uftabduktion $>$120\textdegree{} und erh\"oht das Adduktoren-Verletzungsrisiko signifikant.
  \item \textbf{D6b -- Verletzungsrisiko mit Fremdeinwirkung:} Verletzungen durch Kollisionen mit Gegenspielern, Pfosten oder Boden. Teilweise durch Technikwahl beeinflussbar, da die Technik die Kollisionsgeometrie bestimmt. Beispiel: Ein blockender Torh\"uter im 1v1 liegt am Boden und ist dem Nachtreten des St\"urmers ausgesetzt -- eine Kollisionsgeometrie, die beim Fangen nicht entsteht.
\end{itemize}

$D6(t) = 0.5 \times D6a(t) + 0.5 \times D6b(t)$, invertiert skaliert: 1.0 = niedriges Risiko, 0.0 = hohes Risiko.


\subsection{Anwendung: Fangen vs.\ Blocken}
\label{sec:fangen-vs-blocken}

\begin{table}[H]
\centering
\caption{Fangen vs.\ Blocken im Sechs-Dimensionen-Vergleich}
\label{tab:vergleich}
\begin{tabularx}{\textwidth}{@{}l X X@{}}
\toprule
\textbf{Dimension} & \textbf{Fangen} & \textbf{Blocken} \\
\midrule
D1: Spielrelevanz & Hoch -- Mehrheit der Save-Aktionen & Niedrig -- wenige 1v1-Situationen/Spiel \\[6pt]
D2: Risikoreduktion & Maximal -- Ball gesichert, kein Nachschuss & Mittel -- Ball im Spiel; Nachschuss, Ecke m\"oglich \\[6pt]
D3: Lernaufwand & Hoch, aber kontinuierlich -- stetige Verfeinerung & Hoch -- Timing/K\"orperposition schwer automatisierbar \\[6pt]
D4: Trainerkompetenz & Gut vermittelbar -- etablierte Methodik & Anspruchsvoll -- Entscheidungskontext n\"otig \\[6pt]
D5: Strateg.\ Pot. & Maximal -- Spieler\"offnung, Transition & Minimal -- Mannschaft bleibt defensiv \\[6pt]
\textbf{D6: Verletzungsrisiko} & \textbf{Niedrig (0.70)} -- kontrollierte Bewegungen, Haltung sch\"utzt Gelenke & \textbf{Hoch (0.225)} -- extreme Gelenkwinkel, Kollisionsexposition \\
\bottomrule
\end{tabularx}
\end{table}

\textbf{Fangen ist in allen sechs Dimensionen \"uberlegen.} Mit der Erg\"anzung des Verletzungsrisikos (D6) verliert Blocken auch die letzte Dimension: Die Spread-Technik erfordert extreme H\"uftabduktion ($>$120\textdegree{}) und exponiert den Torh\"uter dem Nachtreten -- das Fangen ist 3.1$\times$ sicherer (D6-Ratio: 0.70 vs.\ 0.225). Blocken hat seine Berechtigung als situatives Werkzeug in spezifischen 1v1-Szenarien -- aber nicht als dominante Trainingsphilosophie.


% =============================================================================
\section{Substitution, Komplementarit\"at und systemische Effekte}
\label{sec:komplementaritaet}
% =============================================================================

\subsection{Substitutionseffekte}
\label{sec:substitution}

\textbf{Zeitliche Substitution:} Jede Trainingsminute f\"ur Blocktechnik fehlt beim Fang-Training, Positionsspiel oder Spieler\"offnung. Eine klassische Opportunit\"atskostenrechnung.

\textbf{Kognitive Substitution:} Der Torh\"uter, der den Block als Default automatisiert, substituiert unbewusst das Fangen in Situationen, wo Fangen besser w\"are -- ein Verf\"ugbarkeitsbias in der Motorik.

\textbf{Systemische Substitution:} Auf Mannschaftsebene substituiert ein blockender Torh\"uter die schnelle Spieler\"offnung durch eine Defensivschlaufe. Die Substitution pflanzt sich durch das gesamte Leistungssystem fort.

\textbf{Verletzungs-Substitution} (neu in v3): Ein verletzter Torh\"uter substituiert Spielzeit durch Rehabilitationszeit. Diese Substitution ist \emph{nicht-linear}: Ein einziger Kreuzbandriss kann eine gesamte Saison kosten. Techniken mit h\"oherem Verletzungsrisiko (Blocken: D6~=~0.225) erh\"ohen die erwartete Ausfallzeit und substituieren damit \emph{Verf\"ugbarkeit} -- die fundamentalste Ressource. Der effektive Wert sinkt: $V^{*}_{\text{eff}} = V^{*} \times A \times D$, wobei $A$ die Verf\"ugbarkeit und $D$ den Dauerhaftigkeitsfaktor beschreibt.


\subsection{Komplementarit\"at der Leistungsebenen}
\label{sec:komplement}

Die drei Leistungsebenen sind untereinander nicht substituierbar, sondern \textbf{komplement\"ar im strengen Sinne}: Der Grenzwert einer Ebene steigt mit dem Niveau der anderen. Operative Exzellenz ohne taktische Entscheidungsqualit\"at ist wertlos. Taktische Intelligenz ohne strategische Einbettung verpufft. Wer nur operativ trainiert, bildet Torh\"uter aus, die technisch sauber blocken, aber nicht wissen, \emph{wann} -- und dem Team strategisch nichts zur\"uckgeben.


\subsection{Das Small-Sample-Problem}
\label{sec:small-sample}

Bei 4--5 Sch\"ussen aufs Tor pro Spiel fehlt die methodische Grundlage, um Noise von Signal zu trennen. Der Social-Media-Effekt verst\"arkt dies: Spektakul\"are Saves haben h\"ohere Verbreitungswahrscheinlichkeit als Routine-Fangaktionen. Ein sauber gefangener Ball geht nicht viral. Ein Spread-Block gegen einen Top-St\"urmer schon. Die Wahrnehmung wird systematisch verzerrt -- ein Mechanismus, der aus der Forschung zur Verf\"ugbarkeitsheuristik bekannt ist \citep{tversky1973availability}.


% =============================================================================
\section{Die Simulator-Analogie: Trainingsdesign nach dem F1-Prinzip}
\label{sec:simulator}
% =============================================================================

Formel-1-Fahrer trainieren im Simulator nicht einzelne Techniken isoliert -- sie trainieren \textbf{Entscheidungen unter realistischen Bedingungen}. Der Simulator bildet die tats\"achliche H\"aufigkeitsverteilung ab: 80\% normale Kurven, 15\% \"Uberholman\"over, 5\% Krisensituationen. Der Fahrer lernt nicht nur \emph{wie}, sondern \emph{wann} welche Reaktion richtig ist.

Das aktuelle Torwarttraining entspricht dem Gegenteil: Block-\"Ubung, Fang-\"Ubung, Parier-\"Ubung als isolierte Stationen. Im Motorsport w\"are das so, als w\"urde ein F1-Fahrer eine Stunde nur Nassfahren \"uben, eine Stunde nur Bremspunkte, eine Stunde nur \"Uberholman\"over -- ohne realistische Abfolge und H\"aufigkeit.


\subsection{Datenbasiertes Trainingsprotokoll}
\label{sec:trainingsprotokoll}

\begin{enumerate}
  \item \textbf{Spielanalyse:} Verteilung der Aktionen des eigenen Torh\"uters \"uber f\"unf Spiele erheben -- Strichliste gen\"ugt.
  \item \textbf{\"Ubungsdesign:} Sequenzen, in denen Situationstypen in realer H\"aufigkeit aufeinanderfolgen -- mit Entscheidungsdruck und ohne Vorank\"undigung.
  \item \textbf{Kalibrierung:} Bewertung jeder Technik \"uber das Sechs-Dimensionen-Framework und Adjustierung der Trainingsallokation.
\end{enumerate}


\subsection{Der Trainer als Szenario-Designer}
\label{sec:szenario-designer}

In diesem Modell wird der Torwarttrainer zum \textbf{Szenario-Designer} -- analog zum F1-Renningenieur. Die Qualit\"at des Trainings h\"angt von der F\"ahigkeit ab, die Spielrealit\"at zu analysieren und \"Ubungen entsprechend zu kalibrieren. Dies erfordert:

\begin{itemize}[nosep]
  \item Diagnostische F\"ahigkeit
  \item Entscheidungsschulung statt reiner Technikschulung
  \item Reflexionsf\"ahigkeit gegen\"uber Trends
\end{itemize}


% =============================================================================
\section{Messbarkeit taktischer und strategischer Entscheidungen}
\label{sec:messbarkeit}
% =============================================================================

\subsection{Was heute messbar ist}
\label{sec:messbar-heute}

Die aktuelle Dateninfrastruktur (Opta, StatsBomb) erfasst diskrete Ereignisse. StatsBomb hat mit GSAA, Positioning Error und CCAA einen wichtigen Schritt geleistet \citep{statsbomb2023}. Aber diese Metriken bleiben auf der operativen Ebene.


\subsection{Was fehlt}
\label{sec:was-fehlt}

Die Erfassung der \textbf{Entscheidungsqualit\"at im Kontext}:

\begin{itemize}[nosep]
  \item Wurde geblockt, obwohl Fangen m\"oglich war?
  \item Hat der Torh\"uter durch Herauslaufen die Schussqualit\"at reduziert?
  \item Hat seine Kommunikation die Abwehrkette in bessere Position gebracht?
\end{itemize}

Dies erfordert die Kombination von Event-Daten, Tracking-Daten und qualitativer Spielanalyse.


\subsection{Forschungsperspektive}
\label{sec:forschung}

Ein vielversprechender Ansatz: Opta Save-Type-Daten (Caught vs.\ Parried Safe vs.\ Parried Danger) mit Spielsituation (1v1, Fernschuss, Standard, Flanke) verkn\"upfen und die resultierenden Anschlussaktionen (Ballbesitz gesichert vs.\ Ecke vs.\ Nachschuss) systematisch erfassen. Dies w\"urde erstmals eine datengest\"utzte Antwort erm\"oglichen auf die Frage, wie oft welche Technik in welcher Situation angewandt wird -- und welche Konsequenzen daraus folgen.


% =============================================================================
\section{Das Visibility Framework: Sichtbarkeit $\neq$ Wertsch\"opfung}
\label{sec:visibility}
% =============================================================================

Eine der zentralen Erkenntnisse des GPM-2.0 ist die systematische Diskrepanz zwischen der \textbf{Sichtbarkeit} einer Aktion und ihrem \textbf{tats\"achlichen Wertsch\"opfungsbeitrag}. Diese Erkenntnis st\"utzt sich auf zwei fundamentale theoretische Grundlagen.

\subsection{Holmstr\"om-Milgrom Multi-Task-Theorie}
\label{sec:holmstrom}

\citet{holmstrom1991} zeigen in ihrer wegweisenden Arbeit zur Multi-Task Principal-Agent-Theorie: Wenn ein Agent mehrere Aufgaben erf\"ullt, die sich in ihrer \textbf{Messbarkeit} unterscheiden, f\"uhren leistungsbasierte Anreize f\"ur die messbaren Aufgaben dazu, dass der Agent seinen Effort systematisch zugunsten der messbaren und zulasten der nicht-messbaren Aufgaben verzerrt. In Extremf\"allen kann ein fixer Lohn optimal sein, um eine ausgewogene Effortverteilung zu gew\"ahrleisten.

Die Anwendung auf die Torwartbewertung ist direkt: Spektakul\"are Saves (messbar, sichtbar) erhalten \"uberproportionale Aufmerksamkeit, w\"ahrend Positionierung, Kommunikation und Defensivorganisation (schwer messbar, unsichtbar) systematisch unterbewertet werden.

\subsection{Kerrs Folly: Rewarding A, While Hoping for B}
\label{sec:kerr}

\citet{kerr1975} dokumentiert das universelle Ph\"anomen misalignierter Anreizsysteme: Organisationen belohnen regelm\"assig Verhaltensweisen, die sie eigentlich verhindern wollen, w\"ahrend die erw\"unschten Verhaltensweisen nicht belohnt werden. Im Torwartkontext: Medien, Fans und teilweise auch Clubs belohnen spektakul\"are Saves (\emph{Rewarding A}), w\"ahrend sie sich eigentlich w\"unschen, dass der Torh\"uter durch gute Positionierung daf\"ur sorgt, dass Sch\"usse gar nicht erst auf sein Tor kommen (\emph{Hoping for B}).

\subsection{Das Visibility-Contribution Spektrum}
\label{sec:vc-spectrum}

Wir formalisieren diese Einsicht als \textbf{Visibility-Contribution Korrelation} $\text{Corr}(V,C)$, die \"uber Positionen hinweg variiert:

\begin{table}[H]
\centering
\caption{Visibility-Contribution Korrelation nach Position}
\label{tab:vc-spectrum}
\begin{tabular}{@{}lrl@{}}
\toprule
\textbf{Position} & $\text{Corr}(V,C)$ & \textbf{Interpretation} \\
\midrule
St\"urmer     & $+0.85$ & Sichtbar = Wertvoll (Tore z\"ahlen direkt) \\
Mittelfeld   & $+0.40$ & Teilweise sichtbar (Assists, Passquote) \\
Verteidiger  & $+0.10$ & Wenig sichtbar (verhinderte Chancen) \\
Torh\"uter    & $-0.15$ & \textbf{Invers} -- beste Leistung = nichts passiert \\
\bottomrule
\end{tabular}
\end{table}

F\"ur den Torh\"uter gilt das \textbf{Prevention Paradox} \citep{rose1985}: Die beste Torh\"uterleistung besteht darin, dass Gegentore gar nicht erst m\"oglich werden -- eine pr\"aventive Leistung, die per Definition unsichtbar bleibt.


% =============================================================================
\section{Counterfactual Performance: Was nicht passiert ist}
\label{sec:counterfactual}
% =============================================================================

Die Sichtbarkeitsproblematik f\"uhrt direkt zur Frage der \textbf{kontrafaktischen Leistung}. Die Gesamtleistung eines Torh\"uters ist:

\begin{equation}
\label{eq:counterfactual}
V^{*} \;=\; V_{\text{beobachtet}} \;+\; V_{\text{verhindert}}
\end{equation}

wobei $V_{\text{beobachtet}}$ die traditionell erfassbaren Aktionen (Saves, GSAA) und $V_{\text{verhindert}}$ die durch Positionierung, Kommunikation und Defensivorganisation vermiedenen Gefahren beschreibt. Unsere LLMMC-Sch\"atzung ergibt, dass $V_{\text{verhindert}} / V^{*} \approx 0.60$ -- die Mehrheit der Torh\"uterleistung ist unsichtbar.

\subsection{State of the Art: PSxG und GSAA}

Die moderne Torwartanalytik hat erste Schritte unternommen, um \"uber reine Save-Quoten hinauszugehen:

\begin{itemize}[nosep]
  \item \textbf{Post-Shot Expected Goals (PSxG):} Ber\"ucksichtigt Schussplatzierung und -geschwindigkeit
  \item \textbf{Goals Saved Above Average (GSAA):} Misst die \"Uber-/Unterperformance relativ zum Durchschnitt
  \item \textbf{Bayesian xGOT:} Probabilistische Modelle mit Prior-Updates
\end{itemize}

Diese Metriken messen die \emph{operative Leistung} besser -- aber die \emph{pr\"aventive Dimension} (Positionierung, Kommunikation, Raumverknappung) bleibt eine offene Forschungsfrage.


% =============================================================================
\section{Action Bias: Warum Torh\"uter zu viel springen}
\label{sec:action-bias}
% =============================================================================

\citet{bareli2007} analysieren 286 Elfmetersituationen und dokumentieren einen systematischen \textbf{Action Bias}: Torh\"uter springen in 93.7\% der F\"alle nach links oder rechts -- obwohl Stehenbleiben in der Mitte die h\"ochste Haltewahrscheinlichkeit bietet ($P(\text{Save}|\text{Mitte}) = 33.3\%$ vs.\ $P(\text{Save}|\text{Sprung}) \approx 25\%$).

Der Mechanismus ist die \textbf{Norm-Theorie} (Kahneman \& Miller, 1986): Das Bedauern \"uber Unterlassung (Stehenbleiben und Tor kassieren) wird als gr\"osser empfunden als das Bedauern \"uber Handeln (Springen und Tor kassieren). Norm-kongruentes Scheitern (Sprung) f\"uhlt sich weniger schlimm an als norm-deviantes Scheitern (Stehenbleiben).

\begin{equation}
\label{eq:action-bias}
\text{Regret}(\text{Unterlassung}) \;>\; \text{Regret}(\text{Handlung}) \quad \text{bei negativem Ergebnis}
\end{equation}

Dieser Befund ist die \textbf{mikroskopische Version} des Visibility-Problems: Auch innerhalb der operativen Ebene bevorzugen Torh\"uter die sichtbare Aktion (Sprung) gegen\"uber der unsichtbaren, aber statistisch \"uberlegenen Strategie (Stehenbleiben).

\subsection{Anwendung auf das Training}

Der Action Bias l\"asst sich durch drei Interventionspunkte adressieren:

\begin{enumerate}[nosep]
  \item \textbf{Statistisches Wissen:} Torh\"uter mit der empirischen Evidenz konfrontieren
  \item \textbf{Normverschiebung:} Stehenbleiben als taktisch kluge Entscheidung umrahmen
  \item \textbf{Wiederholte Exposition:} Im Training Situationen schaffen, in denen Stehenbleiben belohnt wird
\end{enumerate}


% =============================================================================
\section{Generalisierung: Das Visibility-Paradox \"uber Positionen hinweg}
\label{sec:generalisierung}
% =============================================================================

Das Visibility-Paradox des Torh\"uters ist kein Einzelph\"anomen, sondern Ausdruck eines allgemeinen Prinzips: Je pr\"aventiver die Leistung, desto unsichtbarer, und desto systematischer unterbewertet. Tabelle~\ref{tab:vc-spectrum} zeigt das Spektrum von $\text{Corr}(V,C) = +0.85$ (St\"urmer) bis $-0.15$ (Torh\"uter).

Dieses Muster findet sich \"uber den Fussball hinaus: In der Medizin beschreibt \citet{rose1985} das identische Paradox -- pr\"aventive Massnahmen, die der Bev\"olkerung enormen Nutzen bringen, bleiben f\"ur das Individuum unsichtbar. In der Organisationstheorie dokumentiert \citet{kerr1975} dasselbe f\"ur Unternehmen. \citet{holmstrom1991} formalisiert es f\"ur alle Multi-Task-Umgebungen.

\textbf{GPM-2.0 Vorhersage:} Je st\"arker die Corr$(V,C)$ negativ ist, desto gr\"osser die Diskrepanz zwischen traditioneller Bewertung und tats\"achlichem Wert. F\"ur den Fussball bedeutet dies: Die gr\"osste Fehlbewertung findet bei Torh\"utern statt, die kleinste bei St\"urmern.


% =============================================================================
\section{Verletzungsrisiko als Bewertungsdimension}
\label{sec:verletzungsrisiko}
% =============================================================================

Die sechste Dimension D6 schliesst eine fundamentale L\"ucke: Bisherige Leistungsmodelle behandeln Verletzungen als exogenen Schock -- als etwas, das dem Torh\"uter \emph{passiert}. Das GPM-3.0 erkennt, dass Verletzungsrisiko \textbf{endogen} ist: Die Technikwahl bestimmt die Verletzungswahrscheinlichkeit.

\subsection{Verletzungsrisiko ohne Fremdeinwirkung (D6a)}

Non-Contact-Verletzungen entstehen durch die Biomechanik der Technik selbst. Sie sind \textbf{direkt durch Technikwahl und Trainingsdesign beeinflussbar}. Zentrale Risikokategorien:

\begin{itemize}[nosep]
  \item \textbf{Muskel-Sehnen-Verletzungen:} Adduktorenzerrungen bei extremer H\"uftabduktion (Spread-Technik $>$120\textdegree{})
  \item \textbf{Gelenkbelastungen:} Knie-/H\"uftarthrose bei repetitiver Belastung in extremen Winkeln
  \item \textbf{Chronische \"Uberlastung:} Kumulative Sch\"aden durch hohe Trainingsfrequenz in Risikotechniken
\end{itemize}

Fangen erreicht D6a~=~0.85 (niedriges Risiko): Kontrollierte Armbewegungen bei stabiler Rumpfhaltung. Blocken erreicht D6a~=~0.25 (hohes Risiko): Extreme Gelenkwinkel unter dynamischer Belastung.

\subsection{Verletzungsrisiko mit Fremdeinwirkung (D6b)}

Contact-Verletzungen entstehen durch Kollisionen mit Gegenspielern, Pfosten oder Boden. Sie sind \textbf{teilweise durch Technikwahl beeinflussbar}, da die Technik die Kollisionsgeometrie bestimmt:

\begin{itemize}[nosep]
  \item \textbf{Kollisionsgeometrie:} Ein blockender Torh\"uter liegt am Boden und ist dem Nachtreten des St\"urmers ausgesetzt -- eine Geometrie, die beim Fangen nicht entsteht
  \item \textbf{Sch\"utzendes Verhalten:} Stehende/aufrechte Positionen (Fangen) erlauben Ausweichbewegungen; liegende Positionen (Blocken) nicht
  \item \textbf{Kopfverletzungen:} Kopf n\"aher am Boden und an St\"urmerfüssen bei Block-Techniken; erh\"ohtes Konkussionsrisiko
\end{itemize}

Fangen erreicht D6b~=~0.55 (moderates Risiko): Aufrechte Position erm\"oglicht Ausweichen. Blocken erreicht D6b~=~0.20 (hohes Risiko): Maximale Expositionsfl\"ache bei minimaler Ausweichm\"oglichkeit.

\subsection{Verf\"ugbarkeit als Veto-Faktor}

Die zentrale Erkenntnis: \textbf{Ein verletzter Torh\"uter hat einen effektiven Wert von Null.} Der effektive Wert einer Technik muss die Verf\"ugbarkeit ber\"ucksichtigen:

\begin{equation}
\label{eq:veff}
V^{*}_{\text{eff}}(t) = V^{*}(t) \times A(t) \times D(t)
\end{equation}

\noindent wobei $V^{*}$ der Leistungswert, $A \in [0,1]$ die Verf\"ugbarkeit (Anteil der spielbaren Spiele) und $D \in [0,1]$ der Dauerhaftigkeitsfaktor (Leistungserhalt \"uber die Karriere) ist. Verf\"ugbarkeit wirkt als \emph{Veto-Faktor}: $A = 0 \Rightarrow V^{*}_{\text{eff}} = 0$, unabh\"angig von der Leistungsqualit\"at.


% =============================================================================
\section{Praktiker-Validierung: Red Bull Scouting}
\label{sec:redbull}
% =============================================================================

Die theoretischen Ergebnisse des GPM werden durch die Scouting-Praxis von Red Bull gest\"utzt. \citet{wagner2023scouting} dokumentiert das Torh\"uter-Scouting-Profil mit 18 Kompetenzen in f\"unf Kategorien:

\begin{table}[H]
\centering
\caption{Red Bull Scouting-Profil: 18 Kompetenzen nach Wagner (2023)}
\label{tab:redbull}
\begin{tabularx}{\textwidth}{@{}l c X@{}}
\toprule
\textbf{Kategorie} & \textbf{Anzahl} & \textbf{GPM-Mapping} \\
\midrule
Psychologisch & 4 & D1 (Spielrelevanz), D3 (Lernaufwand) \\[4pt]
Sozial & 4 & D4 (Trainerkompetenz), D5 (Strateg.\ Pot.) \\[4pt]
Physisch & 4 & D2 (Risikoreduktion), D6 (Verletzungsrisiko) \\[4pt]
Technisch & 4 & D1, D2, D3 (operativ) \\[4pt]
Taktisch & 2 & D5 (strategisch) \\
\bottomrule
\end{tabularx}
\end{table}

\textbf{Zentrale Erkenntnis:} Von den 18 Kompetenzen sind 8 (44\%) durch reine Datenanalyse nicht messbar -- insbesondere die psychologischen und sozialen Dimensionen. Dies best\"atigt die GPM-These, dass ein \textbf{Hybrid-Modell} aus datengest\"utzter Analyse und Scouting-Expertise erforderlich ist:

\begin{equation}
\label{eq:hybrid}
V^{*}_{\text{hybrid}} = w_{\text{data}} \times V^{*}_{\text{data}} + w_{\text{scout}} \times V^{*}_{\text{scout}}
\end{equation}

\noindent mit $w_{\text{data}} = 0.55$ und $w_{\text{scout}} = 0.45$, basierend auf dem Anteil messbarer vs.\ nicht-messbarer Kompetenzen.


% =============================================================================
\section{Empirische Fallstudien: Drei Torh\"uterprofile}
\label{sec:fallstudien}
% =============================================================================

Das GPM-Framework generiert zw\"olf testbare Vorhersagen (PRED-GPM-001 bis -012). Um diese empirisch zu verankern, evaluieren wir drei Referenz-Torh\"uter, die unterschiedliche Profiltypen repr\"asentieren: den \emph{strategischen Spezialisten}, den \emph{balancierten Allrounder} und den \emph{operativen Spezialisten}.

\subsection{Methodik}

Jeder Torh\"uter wird auf drei Leistungsebenen bewertet ($L_S$, $L_T$, $L_O$), jeweils als Mittelwert aus 4--5 Subkomponenten auf der Skala $[0,1]$. Die Gesamtleistung berechnet sich als:
\[
L(\text{GK}) = w_S \cdot L_S + w_T \cdot L_T + w_O \cdot L_O + \gamma_{ST} \cdot L_S \cdot L_T + \gamma_{SO} \cdot L_S \cdot L_O + \gamma_{TO} \cdot L_T \cdot L_O
\]
mit $w_S = 0.30$, $w_T = 0.35$, $w_O = 0.35$, $\gamma_{ST} = 0.45$, $\gamma_{SO} = 0.30$, $\gamma_{TO} = 0.50$. Die normalisierte Gesamtleistung wird mit $V^{*}$ bezeichnet. Der effektive Wert $V^{*}_{\text{eff}} = V^{*} \times A \times D$ integriert Verf\"ugbarkeit ($A$) und Langlebigkeit ($D$).

Datenquellen: \"Offentliche Matchstatistiken, Scouting-Berichte, Red Bull Kompetenzmodell \citep{wagner2023scouting}, sowie Event-Daten (Opta/StatsBomb).

\subsection{Manuel Neuer: Der unsichtbare Spielmacher}
\label{sec:cs_neuer}

\textbf{Profil:} $L_S = 0.90$, $L_T = 0.75$, $L_O = 0.80$ (Strategischer Spezialist)

Neuer revolution\"arte das Torwartspiel durch Priorisierung der strategischen Ebene: Positionsspiel, Spieler\"offnung, antizipative Abwehrorganisation. Sein $L_S = 0.90$ ist der h\"ochste Wert im Referenzpanel. Trotz altersbedingtem R\"uckgang der operativen Werte ($L_O = 0.80$) bleibt seine Gesamtleistung dank Komplementarit\"at hoch:
\[
L(\text{Neuer}) = 0.30 \times 0.90 + 0.35 \times 0.75 + 0.35 \times 0.80 + \gamma\text{-Terme} = 0.725 \quad (\text{normalisiert})
\]

\textbf{Prediction-Validierung:}
\begin{itemize}
\item \textbf{PRED-006} (Positioning $>$ Spectacle): Best\"atigt. Neuers beste Spiele haben die wenigsten Paraden -- das \emph{Prevention Paradox} in Reinform.
\item \textbf{PRED-008} (Pr\"aventionsparadox): Best\"atigt. Seine niedrigste Paraden-pro-Spiel-Rate korreliert mit den besten Defensivleistungen.
\end{itemize}

\textbf{Trainingsallokation:} Fangen 38\%, Abwehr 22\%, Blocken 5\%, Fausten 15\%, Fusssave 8\%, Spielaufbau 12\%. Die niedrige Block-Quote (5\%) entspricht dem GPM-Optimum.

\textbf{V*$_{\text{eff}}$-Analyse:} In der Saison 2023-24 (nach Beinbruch) sank $A = 0.64$, was zu $V^{*}_{\text{eff}} = 0.725 \times 0.64 = 0.464$ f\"uhrte -- ein signifikanter, aber nicht katastrophaler R\"uckgang dank hoher Grundqualit\"at.

\subsection{Alisson Becker: Komplementarit\"at in Aktion}
\label{sec:cs_alisson}

\textbf{Profil:} $L_S = 0.80$, $L_T = 0.80$, $L_O = 0.85$ (Balanced Excellence)

Alisson repr\"asentiert den balancierten Profiltyp: keine extremen Spitzen, keine Schw\"achen. Dieser Profiltyp maximiert die Komplementarit\"atsterme, weil $\gamma_{ij} \cdot L_i \cdot L_j$ bei zwei hohen Faktoren gr\"osser ist als bei einem sehr hohen und einem niedrigen:
\[
\gamma\text{-Terme}(\text{Alisson}) = 0.45 \times 0.80^2 + 0.30 \times 0.80 \times 0.85 + 0.50 \times 0.80 \times 0.85 = 0.832
\]

\textbf{Prediction-Validierung:}
\begin{itemize}
\item \textbf{PRED-004} (Balanced $\gamma$-Bonus): Best\"atigt. Alissons Komplementarit\"atsanteil (50.4\%) ist der h\"ochste aller drei Profile. Der ausbalancierte Torh\"uter profitiert \"uberproportional vom Zusammenspiel der Ebenen.
\item \textbf{PRED-002} (Optimale Allokation): Best\"atigt. Seine Trainingsverteilung ($\sim$35\% Fangen, $\sim$6\% Blocken) liegt nahe am GPM-Optimum.
\end{itemize}

\textbf{V*$_{\text{eff}}$-Analyse:} Mit $V^{*} = 0.733$ und $A = 0.88$ (hohe Verf\"ugbarkeit) erreicht Alisson den h\"ochsten effektiven Wert: $V^{*}_{\text{eff}} = 0.649$ -- der beste Wert im gesamten Panel.

\subsection{Thibaut Courtois: Das Verf\"ugbarkeits-Veto}
\label{sec:cs_courtois}

\textbf{Profil:} $L_S = 0.55$, $L_T = 0.65$, $L_O = 0.88$ (Operativer Spezialist)

Courtois ist der paradigmatische Fall f\"ur $V^{*}_{\text{eff}} = V^{*} \times A \times D$. Als operativer Spezialist ($L_O = 0.88$, h\"ochster Einzelwert) verf\"ugt er \"uber beeindruckende physische Attribute: Spannweite, Reflexe, Pr\"asenz im Strafraum. Doch sein Kreuzbandriss im August 2023 reduzierte die Verf\"ugbarkeit auf $A = 0.15$ (8 von 52 Spielen):
\[
V^{*}_{\text{eff}} = 0.574 \times 0.15 \times 1.00 = 0.088
\]

\textbf{Prediction-Validierung:}
\begin{itemize}
\item \textbf{PRED-011} (V*$_{\text{eff}}$-Kollaps): Best\"atigt. 83\% des potenziellen Wertes ($0.517 - 0.088 = 0.429$) wurden durch die Verletzung zerst\"ort.
\item \textbf{PRED-009} (Block-heavy $\rightarrow$ h\"ohere Verletzungsrate): Konsistent. Courtois alloziert 12\% der Trainingszeit auf Blocken (vs. 5--6\% bei Neuer/Alisson). Blocken ist mit $D6 = 0.225$ die gef\"ahrlichste Technik ($3.1\times$ riskanter als Fangen).
\end{itemize}

\textbf{Kontrastierung:} Ein verf\"ugbarer Alisson ($V^{*}_{\text{eff}} = 0.649$) \"ubertrifft Courtois ($V^{*}_{\text{eff}} = 0.088$) um den Faktor 7.4 -- obwohl Courtois h\"ohere operative Rohwerte aufweist.

\subsection{Vergleichende Analyse}
\label{sec:cs_vergleich}

\begin{table}[h]
\centering
\small
\begin{tabular}{lcccccccc}
\hline
\textbf{Torh\"uter} & $L_S$ & $L_T$ & $L_O$ & $V^{*}$ & $\gamma$-\textbf{Anteil} & $A$ & $V^{*}_{\text{eff}}$ & \textbf{Typ} \\
\hline
Neuer     & 0.90 & 0.75 & 0.80 & 0.725 & 50.2\% & 0.64 & 0.464 & Strategisch \\
Alisson   & 0.80 & 0.80 & 0.85 & 0.733 & 50.4\% & 0.88 & 0.649 & Balanced \\
Courtois  & 0.55 & 0.65 & 0.88 & 0.574 & 45.8\% & 0.15 & 0.088 & Operativ \\
\hline
\end{tabular}
\caption{GPM-3.0 Referenzprofile: Drei-Ebenen-Bewertung, Komplementarit\"at und effektiver Wert.}
\label{tab:referenzprofile}
\end{table}

Drei zentrale Befunde ergeben sich aus dem Vergleich:

\textbf{Erstens:} Balance schl\"agt Spezialisierung. Alisson erzielt den h\"ochsten $V^{*}$ trotz niedrigerer Maximalwerte, weil die Komplementarit\"atsterme bei balancierten Profilen maximal werden ($\gamma_{ij} \cdot L_i \cdot L_j$ ist bei $(0.80, 0.80)$ gr\"osser als bei $(0.55, 0.90)$).

\textbf{Zweitens:} Verf\"ugbarkeit ist ein Veto-Faktor. Courtois' operative \"Uberlegenheit ($L_O = 0.88$) wird durch $A = 0.15$ neutralisiert. Die multiplikative Struktur $V^{*}_{\text{eff}} = V^{*} \times A$ bestraft Ausf\"alle exponentiell.

\textbf{Drittens:} Strategische Investition zahlt sich langfristig aus. Neuers hoher $L_S$-Wert kompensiert operative Alterungseffekte und generiert Wert, der von traditionellen Metriken (Paraden/Spiel) nicht erfasst wird.


% =============================================================================
\section{Zusammenfassung und Ausblick}
\label{sec:zusammenfassung}
% =============================================================================

\textbf{Erstens:} Die Torwartleistung setzt sich aus strategischer, taktischer und operativer Leistung zusammen -- jeweils im Leistungsverbund mit der Mannschaft.

\textbf{Zweitens:} Jede Technik l\"asst sich anhand von sechs Dimensionen bewerten. Fangen ist in allen sechs Dimensionen \"uberlegen.

\textbf{Drittens:} Zwischen Techniken und Leistungsebenen bestehen hochgradige Substitutions- und Komplementarit\"atseffekte, die bei isoliertem Training \"ubersehen werden.

\textbf{Viertens:} Die aktuelle Trainingsallokation wird durch einen Social-Media-verst\"arkten Verf\"ugbarkeitsbias verzerrt.

\textbf{F\"unftens:} Ein datenbasiertes Trainingsprotokoll nach dem F1-Simulatorprinzip kann diese Verzerrung korrigieren.

\textbf{Sechstens} (neu in v2): Die Sichtbarkeit einer Aktion korreliert negativ mit ihrem pr\"aventiven Wert. F\"ur Torh\"uter gilt $\text{Corr}(V,C) < 0$ -- die beste Leistung ist unsichtbar (\emph{Prevention Paradox}).

\textbf{Siebtens} (neu in v2): Die Gesamtleistung $V^{*} = V_{\text{beobachtet}} + V_{\text{verhindert}}$ zeigt, dass ca.\ 60\% des Torh\"uterwerts aus pr\"aventiven, unsichtbaren Aktionen stammt.

\textbf{Achtens} (neu in v2): Der Action Bias (93.7\% Sprungquote bei Elfmetern) ist die Mikro-Manifestation des Visibility-Paradoxes und kann durch gezieltes Training adressiert werden.

\textbf{Neuntens} (neu in v3): Verletzungsrisiko ist keine exogene Variable, sondern \textbf{endogen durch Technikwahl bestimmt}. Blocken ist mit D6~=~0.225 die verletzungstr\"achtigste Technik -- Fangen ist 3.1$\times$ sicherer (D6~=~0.70). Die Unterscheidung in D6a (ohne Fremdeinwirkung) und D6b (mit Fremdeinwirkung) zeigt, dass Blocken in \emph{beiden} Kategorien das h\"ochste Risiko aufweist.

\textbf{Zehntens} (neu in v3): Der effektive Wert einer Technik muss Verf\"ugbarkeit als Veto-Faktor ber\"ucksichtigen: $V^{*}_{\text{eff}} = V^{*} \times A \times D$. Ein verletzter Torh\"uter hat einen effektiven Wert von Null, unabh\"angig von seiner Leistungsqualit\"at.

\textbf{Elftens} (neu in v3): Die Validierung durch das Red Bull Scouting-Profil \citep{wagner2023scouting} zeigt, dass 44\% der 18 Torh\"uter-Kompetenzen durch Daten allein nicht messbar sind. Ein Hybrid-Modell ($V^{*}_{\text{hybrid}} = 0.55 \times V^{*}_{\text{data}} + 0.45 \times V^{*}_{\text{scout}}$) ist erforderlich.

\textbf{Zw\"olftens} (neu in v3): Die empirischen Fallstudien (Neuer, Alisson, Courtois) validieren die zentralen Predictions: Balance schl\"agt Spezialisierung ($\gamma$-Bonus), Verf\"ugbarkeit ist ein Veto-Faktor ($V^{*}_{\text{eff}}$-Kollaps), und strategische Investition generiert unsichtbaren, aber messbaren Wert.

\vspace{1em}
\noindent Dieses Framework versteht sich als Einladung an die Torwarttrainer-Community und die Sportwissenschaft, die Diskussion \"uber Trainingsphilosophien auf eine evidenzbasierte Grundlage zu stellen. Die Integration der Multi-Task-Theorie (Holmstr\"om \& Milgrom), des Prevention Paradox (Rose), der Action-Bias-Forschung (Bar-Eli et~al.), der Verletzungsepidemiologie (Ekstrand et~al.) und der Scouting-Praxis (Wagner/Red Bull) bietet erstmals eine theoretisch fundierte Erkl\"arung daf\"ur, \emph{warum} die Torwartbewertung systematisch verzerrt ist -- und \emph{wie} Trainingsprotokolle diese Verzerrung korrigieren k\"onnen.


% =============================================================================
% REFERENCES
% =============================================================================
\newpage
\bibliographystyle{apalike}

\begin{thebibliography}{99}

\bibitem[Bar-Eli et~al., 2007]{bareli2007}
Bar-Eli, M., Azar, O.\,H., Ritov, I., Keidar-Levin, Y. \& Schein, G. (2007).
\newblock Action bias among elite soccer goalkeepers: The case of penalty kicks.
\newblock \emph{Journal of Economic Psychology}, 28(5), 606--621.

\bibitem[Chiappori et~al., 2002]{chiappori2002}
Chiappori, P.\,A., Levitt, S. \& Groseclose, T. (2002).
\newblock Testing mixed-strategy equilibria when players are heterogeneous: The case of penalty kicks in soccer.
\newblock \emph{American Economic Review}, 92(4), 1138--1151.

\bibitem[Gramage Medina et~al., 2025]{gramage2025}
Gramage Medina, K., et~al. (2025).
\newblock Coaches' perceptions on the Spread: A 1 vs 1 goalkeeper technique.
\newblock \emph{International Journal of Sports Science \& Coaching}.

\bibitem[Otte et~al., 2022]{otte2022}
Otte, F.\,W., Dittmer, T. \& West, J. (2022).
\newblock Goalkeeping in modern football: Current positional demands and research insights.
\newblock \emph{International Sport Coaching Journal}, 10, 112--120.

\bibitem[Soccerment, 2024]{soccerment2024}
Soccerment (2024).
\newblock The Vitruvian Goalie: A deep dive into the goalkeepers' stats.
\newblock \emph{soccerment.com}.

\bibitem[Stats Perform / Opta, 2024]{opta2024}
Stats Perform / Opta (2024).
\newblock Opta Event Definitions.
\newblock \emph{statsperform.com}.

\bibitem[StatsBomb, 2023]{statsbomb2023}
StatsBomb (2023).
\newblock Introducing Goalkeeper Radars.
\newblock \emph{blogarchive.statsbomb.com}.

\bibitem[Tigges, 2026]{tigges2026}
Tigges (2026).
\newblock [Reference to tactical and strategic goalkeeper dimensions].

\bibitem[Holmstr\"om \& Milgrom, 1991]{holmstrom1991}
Holmstr\"om, B. \& Milgrom, P. (1991).
\newblock Multitask principal-agent analyses: Incentive contracts, asset ownership, and job design.
\newblock \emph{Journal of Law, Economics, and Organization}, 7, 24--52.

\bibitem[Kerr, 1975]{kerr1975}
Kerr, S. (1975).
\newblock On the folly of rewarding A, while hoping for B.
\newblock \emph{Academy of Management Journal}, 18(4), 769--783.

\bibitem[Palacios-Huerta, 2003]{palacioshuerta2003}
Palacios-Huerta, I. (2003).
\newblock Professionals play minimax.
\newblock \emph{Review of Economic Studies}, 70(2), 395--415.

\bibitem[Rose, 1985]{rose1985}
Rose, G. (1985).
\newblock Sick individuals and sick populations.
\newblock \emph{International Journal of Epidemiology}, 14(1), 32--38.

\bibitem[Tversky \& Kahneman, 1973]{tversky1973availability}
Tversky, A. \& Kahneman, D. (1973).
\newblock Availability: A heuristic for judging frequency and probability.
\newblock \emph{Cognitive Psychology}, 5(2), 207--232.

\bibitem[Wagner, 2023]{wagner2023scouting}
Wagner, H. (2023).
\newblock Scouting und Talentf\"orderung im modernen Torwartspiel: Das Red Bull Kompetenzmodell.
\newblock \emph{Red Bull Football Academy Internal Report}.

\bibitem[Ekstrand et~al., 2011]{ekstrand2011injuries}
Ekstrand, J., H\"agglund, M. \& Wald\'en, M. (2011).
\newblock Epidemiology of muscle injuries in professional football (soccer).
\newblock \emph{American Journal of Sports Medicine}, 39(6), 1226--1232.

\bibitem[Drawer \& Fuller, 2002]{drawer2002goalkeeper}
Drawer, S. \& Fuller, C.\,W. (2002).
\newblock An economic framework for assessing the impact of injuries in professional football.
\newblock \emph{Safety Science}, 40(6), 537--556.

\end{thebibliography}


\end{document}
